\documentclass[11pt]{amsart}
\usepackage{amssymb,amsmath}
\newtheorem{theorem}{Theorem}[section]
\newtheorem{lemma}[theorem]{Lemma}
\newtheorem{proposition}[theorem]{Proposition}
\theoremstyle{remark}\newtheorem{remark}[theorem]{Remark}
\newcommand{\B}{\mathbb B}
\title[Adjoining a zero to a commutative ring]{The semiring obtained by adjoining a zero to a commutative ring:\\
spectrum, modules and projectives}
\author[AI-produced note]{AI-produced note (draft, 21 November 2026)}
\begin{document}
\begin{abstract}
For a nonzero commutative ring $R$ let $G(R)=R\sqcup\{\tau\}$, with $\tau$ adjoined as additive identity and
multiplicative absorber. We describe $\operatorname{Spec}G(R)$ as $\operatorname{Spec}R$ with one Boolean generic
point, show that $G(R)$-semimodules are exactly commutative Clifford monoids with $R$-linear fibres (support diagrams
over join-semilattices), and classify the finitely generated projective $G(R)$-semimodules: their skeleta are the
finite distributive lattices and their fibres are sums of projective $R$-modules indexed by order ideals.
\end{abstract}
\maketitle

\section{The construction}
Let $R\neq0$ be a commutative ring and $G(R)=R\sqcup\{\tau\}$ with $\tau+x=x$, $\tau x=\tau$ for all $x$, and the
operations of $R$ on $R$. Write $e=0_R$; then $e\neq\tau$, $e+e=e$ and $e^2=e$. The forgetful functor $U$ from
commutative semirings to commutative semirings without distinguished zero has the left adjoint $A\mapsto A\sqcup\{\tau\}$
(adjoining a zero, \cite{golan}), and $G=$ this left adjoint composed with $U$. Let $\B=\{0,1\}$, $1+1=1$.

\begin{proposition}\label{normal}
$x\mapsto(p(x),\chi(x))$, with $p(\tau)=0$, $p|_R=\mathrm{id}$, $\chi(\tau)=0$, $\chi|_R=1$, is an injective semiring map
$G(R)\to R\times\B$ with image $\{(x,b):b=0\Rightarrow x=0\}$. Every semiring map from $G(R)$ to a ring factors
through $p$, and $\chi$ is the only semiring map $G(R)\to\B$.
\end{proposition}
\begin{proof}
The operations are coordinatewise on the image and the map is injective. If $f\colon G(R)\to A$ with $A$ a ring, then
$f(e)+f(e)=f(e)$ forces $f(e)=0$, so $f$ factors through $p$. A unital map $g\colon G(R)\to\B$ sends the unit $-1$ to $1$, so
$g(e)=g(1)+g(-1)=1$; then for $x\in R$, $g(x)=g(x)g(e)=g(xe)=g(e)=1$, and $g(\tau)=0$. So $g=\chi$.
\end{proof}

\section{The spectrum}
\begin{theorem}\label{spec}
The ideals of $G(R)$ are $\{\tau\}$ and $J\cup\{\tau\}$ for ideals $J$ of $R$; its prime ideals are $\{\tau\}$ and
$P\cup\{\tau\}$ for $P\in\operatorname{Spec}R$. The point $\eta=\{\tau\}$ lies in every prime, $\{\eta\}$ is open, the
closed set $V(e)$ is homeomorphic to $\operatorname{Spec}R$ through $p$, and the localization at $\eta$ is $\B$. If
$\dim R<\infty$ then $\dim G(R)=\dim R+1$.
\end{theorem}
\begin{proof}
An ideal $I\neq\{\tau\}$ contains some $r\in R$, hence $e=0\cdot r$; then $I\cap R$ is an ideal of $R$ and
$I=(I\cap R)\cup\{\tau\}$. Primality transfers because $xy=\tau$ only if $x=\tau$ or $y=\tau$. Every prime other than
$\eta$ contains $e$, so $\{\eta\}$ is the complement of $V(e)$. Localizing at $\eta$ inverts every element of $R$,
including the idempotent $e$, which forces $e=1$ and identifies all of $R$ with $1$, leaving $\{\tau,1\}\cong\B$.
A chain of primes of $R$ of length $d$ gives one of length $d+1$ in $G(R)$ by adjoining $\eta$ below it.
\end{proof}

\begin{proposition}\label{glue}
If $R$ is a domain with fraction field $K$, then $G(R)\cong R\times_K G(K)$, so that $\operatorname{Spec}G(R)$ is
obtained by gluing the two-point space $\operatorname{Spec}G(K)$ along its closed point to the generic point of
$\operatorname{Spec}R$. For nonzero rings $A,B$, $G(A\times B)\cong G(A)\times_{\B}G(B)$ over the maps $\chi$.
\end{proposition}
\begin{proof}
The elements of $R\times_K G(K)$ are $(r,r)$ for $r\in R$ and $(0,\tau)$, with coordinatewise operations matching
those of $G(R)$. The elements of $G(A)\times_\B G(B)$ are $(\tau,\tau)$ and the pairs in $A\times B$; send $(a,b)$ to
$(a,b)\in A\times B\subset G(A\times B)$ and $(\tau,\tau)$ to $\tau$.
\end{proof}

\section{Modules}
A $G(R)$-semimodule is a commutative monoid $(M,+,0_M)$ with an action of $G(R)$ that is distributive in both
variables, with $1m=m$ and $\tau m=0_M$. A \emph{support diagram} is a join-semilattice $L$ with least element, an
$R$-module $A_l$ for each $l\in L$ and $R$-linear maps $\rho_{l,l'}\colon A_l\to A_{l'}$ for $l\le l'$ forming a functor.

\begin{theorem}\label{modules}
For a $G(R)$-semimodule $M$ put $\varepsilon(m)=em$, $L_M=eM$ and $M_l=\varepsilon^{-1}(l)$. Then $(M,+)$ is a
commutative Clifford monoid: each $M_l$ is an abelian group with identity $l$, and $L_M$ is a join-semilattice under
$+$. Each $M_l$ is an $R$-module, the maps $T_{l,l'}(m)=m+l'$ form a support diagram, and
$m+n=T_{\varepsilon(m),j}(m)+T_{\varepsilon(n),j}(n)$ in $M_j$, $j=\varepsilon(m)+\varepsilon(n)$. This gives an
equivalence between $G(R)$-semimodules and support diagrams; cancellative semimodules correspond to $R$-modules.
\end{theorem}
\begin{proof}
$m+\varepsilon(m)=(1+e)m=m$ and $m+(-1)m=(1+(-1))m=\varepsilon(m)$, and $\varepsilon(m)+\varepsilon(m)=(e+e)m=
\varepsilon(m)$; so $M$ is a union of the groups $M_l$, and $\varepsilon$ is additive. This is Clifford's description of
commutative inverse monoids as strong semilattices of groups \cite{clifford,howie}, with structure maps $T_{l,l'}$.
Since $r\varepsilon(m)=\varepsilon(rm)$, the ring $R$ acts on each fibre, linearly, and the $T_{l,l'}$ are $R$-linear.
Conversely a support diagram defines $\bigsqcup_l A_l$ with $x\boxplus y=\rho(x)+\rho(y)$ in the fibre over the join and
$\tau x=0\in A_{0}$; the two constructions are mutually inverse. A Clifford monoid is cancellative iff it has one
idempotent, that is, iff $L_M$ is a point.
\end{proof}

\section{Projectives}
\begin{lemma}\label{retract}
A finite join-semilattice with least element is a retract of a finite Boolean lattice $\mathcal P([n])$ (by maps
preserving joins and the least element) if and only if it is a distributive lattice.
\end{lemma}
\begin{proof}
If $L$ is such a retract via $i,r$, then $\epsilon=i\circ r$ is an idempotent join-endomorphism of $\mathcal P([n])$
and $L\cong\epsilon(\mathcal P([n]))$ as posets. The image is closed under unions, and its meet is
$U\wedge V=\epsilon(U\cap V)$. As $\epsilon$ preserves unions,
$U\wedge(V\cup W)=\epsilon((U\cap V)\cup(U\cap W))=(U\wedge V)\cup(U\wedge W)$. Conversely, by Birkhoff's theorem a
finite distributive lattice is the lattice $\operatorname{Idl}(J)$ of order ideals of a finite poset $J$, a retract of
$\mathcal P(J)$ through $S\mapsto{\downarrow}S$.
\end{proof}

For a finite poset $J$ and f.g. projective $R$-modules $Q_j$, let $P(J,Q)$ be the support diagram on
$\operatorname{Idl}(J)$ with fibre $\bigoplus_{j\in U}Q_j$ over $U$ and the summand inclusions as transports.

\begin{theorem}\label{proj}
A finitely generated $G(R)$-semimodule is projective if and only if it is isomorphic to some $P(J,Q)$.
\end{theorem}
\begin{proof}
\emph{If.} $P(J,R)=\{x\in G(R)^J:\operatorname{supp}x\in\operatorname{Idl}(J)\}$ is a subsemimodule, and the map that
keeps $x$ on its support, puts the supported zero $e$ on ${\downarrow}\operatorname{supp}x\setminus\operatorname{supp}x$
and $\tau$ elsewhere is a semimodule retraction $G(R)^J\to P(J,R)$. For free fibres $R^{n_j}$ the same map on
$G(R)^{J'}$, $J'=\bigsqcup_j\{j\}\times[n_j]$, with the union-preserving closure $S\mapsto{\downarrow}\pi(S)\times[n_\ast]$
in place of ${\downarrow}$, is a retraction onto $P(J,(R^{n_j}))$. Fibrewise projections then exhibit $P(J,Q)$ as a
retract of $P(J,Q\oplus Q')\cong P(J,(R^{n_j}))$ with $Q_j\oplus Q'_j\cong R^{n_j}$.

\emph{Only if.} Let $M$ be a retract of $F=G(R)^n$ via $i,r$. They commute with $\varepsilon$, so they are maps of
support diagrams. By Lemma~\ref{retract}, $eM\cong\operatorname{Idl}(J)$. Each $M_l$ is a retract of
$R^{i_0(l)}$, hence f.g. projective. Each transport is split injective, with left inverse $r\circ\pi\circ i$ for a
coordinate projection $\pi$. Writing $l\wedge l'=r_0(i_0(l)\cap i_0(l'))$, transport through $i$ and $r$ gives
$M_{l\vee l'}=T(M_l)+T(M_{l'})$ and $T(M_l)\cap T(M_{l'})=T(M_{l\wedge l'})$; for the second, lift $i(x)$ to the fibre
over $i_0(l)\cap i_0(l')$ by replacing its supported zeros outside that set by $\tau$. Thus $l\mapsto M_l$ is a lattice
map into the direct summands of the top fibre. Put $Q_j=M_{\downarrow j}/M_{\downarrow j\setminus\{j\}}$, which is
projective, and choose sections $s_j$. Along a linear extension of $J$, the modular isomorphism shows that
$(q_j)\mapsto\sum s_j(q_j)$ is an isomorphism $\bigoplus_{j\in U}Q_j\to M_U$ compatible with inclusions.
\end{proof}

\begin{remark}
Lemma~\ref{retract} is the finite case of the description of projective semilattices in \cite{hk} (compare
\cite{bh}); the proof above is self-contained, with Birkhoff's representation \cite{birkhoff}.
\end{remark}

\begin{remark}
The retraction proving that $P(J,R)$ is projective fills the down-closure of a support with supported zeros $e$: the
classification rests on the distinction between $e$ and $\tau$.
\end{remark}

\section{Quotients and counts}
The congruences of $G(R)$ are, for each ideal $J$ of $R$, the one collapsing $\tau$ into $J$ (quotient $R/J$) and the one
keeping $\tau$ apart (quotient $G(R/J)$, the based set $(R/J)_+$). For $R=\mathcal O_K$ the first family counts
$\sum_J N(J)^{-s}=\zeta_K(s)$. The second counts $\sum_J(N(J)+1)^{-s}$, whose weight $w=N+1$ satisfies
$w(JK)-1=(w(J)-1)(w(K)-1)$ for coprime $J,K$: by Proposition~\ref{glue} the based quotient of a product is the smash
product of the based quotients. For $R=\mathbb Z$ the two families together count $2\zeta(s)-1$.

\medskip\noindent\emph{Provenance.} The statements are taken from AI-produced research notes (ChatGPT, GPT-5.x Pro,
2026); this note's proofs and checks are by Claude Opus 5.5 at effort max (Anthropic), with exact computations of
Theorem~\ref{spec} for $G(\mathbb Z/n)$, of Proposition~\ref{glue} for $\mathbb Z/6$, of Theorem~\ref{modules} on 80
random subsemimodules of $G(\mathbb F_p)^n$, and of Lemma~\ref{retract} for $n\le4$.

\begin{thebibliography}{9}
\bibitem{clifford} A. H. Clifford, Semigroups admitting relative inverses, Ann. of Math. 42 (1941), 1037--1049.
\bibitem{golan} J. S. Golan, Semirings and their Applications, Kluwer, 1999.
\bibitem{birkhoff} G. Birkhoff, Rings of sets, Duke Math. J. 3 (1937), 443--454.
\bibitem{hk} A. Horn and N. Kimura, The category of semilattices, Algebra Universalis 1 (1971), 26--38.
\bibitem{bh} R. Balbes and A. Horn, Projective distributive lattices, Pacific J. Math. 33 (1970), 273--279.
\bibitem{howie} J. M. Howie, Fundamentals of Semigroup Theory, Oxford University Press, 1995.
\end{thebibliography}
\end{document}
